%%%%%%%%%%%%%%%%%%%%%%%%%%%%%%%%%%%%%%%%%%%%%%%%%%%%%%%%%%%%%%%%%%%%%%%%%%%%%%%%%%%%%%%%%%%%%%%%%%%%%%
%	COMMENTS
%%%%%%%%%%%%%%%%%%%%%%%%%%%%%%%%%%%%%%%%%%%%%%%%%%%%%%%%%%%%%%%%%%%%%%%%%%%%%%%%%%%%%%%%%%%%%%%%%%%%%%
% Headline is created automatically

As part of this project, an eye tracker was implemented in a driving simulator at Karlsruhe University of Applied Sciences using inexpensive, commercially available hardware. The aim was to record the gaze position of test subjects in various traffic scenarios in order to analyse their perception and decision-making processes in road traffic more systematically. The aim of the project was not to develop a completely new tracking system from scratch. Instead, existing eye tracking systems were to be evaluated for their applicability in the university's driving simulator and implemented. The main part of the project consisted of finding a suitable eye tracking system, implementing it in the CARLA environment and finding solutions to any technical problems that might arise. 

After a market analysis, the Beam Eye Tracker from Eyeware Tech was selected due to its user-friendly API, low purchase price, low hardware requirements and accuracy that met the requirements of the application. The implementation was carried out using Python with the Beam SDK. The tracking data was linked to CARLA, in particular to the data from the Instance Segmentation Sensor, in order to infer the objects viewed in the simulation from the gaze position estimated by the eye tracker. This allows the eye tracker not only to output the test person's gaze position in x and y coordinates relative to the screen, but also to recognise which object the test person is looking at within the CARLA simulation environment at any given time. The interface between the eye tracker data and the CARLA Instance Segmentation Sensor data was tested in a specially created test environment in CARLA. 

A key technical problem was the curved projection screen of the driving simulator, which led to distortions in the recorded gaze data. To compensate for this, a software-based correction was implemented that uses mathematical functions to minimise the error caused by the curved screen. The accuracy of the implemented tracking system is sufficient to be able to determine which object the test person is looking at at any given time for larger objects in the CARLA simulation world. These are good results, especially considering the low financial outlay compared to eye tracking systems that cost several hundred euros. It was possible to implement the eye tracking system free of charge. The only costs incurred were those of the webcam used. Implementation in the driving simulator and in the CARLA environment was completed in full.

The work proves that even with inexpensive hardware and software, it is possible to functionally integrate eye tracking systems into driving simulators with large, curved screens, thus laying the foundation for future experiments in driving behaviour analysis. The work also identifies potential approaches for further improving tracking accuracy in the future. 